\documentclass[11pt]{article}
\usepackage[a4paper,margin=1in]{geometry}
\usepackage{amsmath,amssymb,amsthm,mathtools}
\usepackage{booktabs}
\usepackage{tabularx}
\usepackage{array}
\usepackage{enumitem}
\usepackage[dvipsnames]{xcolor}
\usepackage{hyperref}
\usepackage{microtype}
\usepackage{comment}

\hypersetup{
  colorlinks=true,
  linkcolor=MidnightBlue,
  urlcolor=MidnightBlue
}

% Define \WITHSOLUTIONS before inputting this file to show solutions.
\newif\ifsolutions
\ifdefined\WITHSOLUTIONS
  \solutionstrue
\else
  \solutionsfalse
\fi

\newcommand{\Pbb}{\mathbb{P}}
\newcommand{\E}{\mathbb{E}}
\newcommand{\rv}[1]{\widetilde{#1}}
\newcommand{\po}[1]{\widetilde{\mathrm{po}}_{#1}}
\newcommand{\indep}{\mathrel{\perp\!\!\!\perp}}
\newcommand{\given}{\,\middle|\,}

\ifsolutions
  \newenvironment{answer}{%
    \par\medskip\noindent\textcolor{MidnightBlue}{\textbf{Solution.}}\ }
    {\par\medskip}
  \newcommand{\workingspace}[1]{}
\else
  \excludecomment{answer}
  \newcommand{\workingspace}[1]{\vspace{#1}}
\fi

\newcommand{\keybox}[1]{%
  \begin{center}
  \fcolorbox{MidnightBlue}{MidnightBlue!5}{%
    \begin{minipage}{0.91\textwidth}#1\end{minipage}}
  \end{center}}

\title{\huge Recitation 2\\[0.35em]
  \normalsize Causal Inference, Confounding, and Simpson's Paradox\\[0.35em]
  \large\ifsolutions Solutions and Teaching Notes\else Review and Practice Problems\fi}
\author{Xuanxi Zhang}
\date{September 18, 2026}

\begin{document}
\maketitle

\ifsolutions
\section*{Pacing guide (75 minutes)}
\begin{center}
\begin{tabular}{@{}lll@{}}
\toprule
Time & Activity & Suggested use \\
\midrule
0--25 min & Review & Potential outcomes, confounding, Curry--Lee example \\
25--38 min & Problem 1 & 7 min individual work, then discussion \\
38--53 min & Problem 2 & 8 min individual work, then discussion \\
53--68 min & Problem 3 & 8 min individual work, then discussion \\
68--75 min & Wrap-up or Problem 4 & Use Problem 4 only if the class is ahead \\
\bottomrule
\end{tabular}
\end{center}

\paragraph{Teaching emphasis.}
Keep three quantities distinct: an observed conditional probability, a
potential-outcome probability, and a standardized probability. In the
basketball example, emphasize that standardization is a better descriptive
comparison but is causal only under the stated conditional-independence
assumption.
\fi

\section{Potential outcomes and causal effects}

Let the binary treatment $\rv{t}$ equal $1$ for treatment and $0$ for control.
Every unit has two \textbf{potential outcomes}:
\[
  \po{0}=\text{outcome under control},
  \qquad
  \po{1}=\text{outcome under treatment}.
\]
Both are defined whether or not the unit actually receives treatment. The
observed outcome is
\[
  \rv{y}=
  \begin{cases}
    \po{0}, & \rv{t}=0,\\
    \po{1}, & \rv{t}=1,
  \end{cases}
  \qquad\text{or equivalently}\qquad
  \rv{y}=(1-\rv{t})\po{0}+\rv{t}\po{1}.
\]
This relationship is called \textbf{consistency}. For one unit we observe only
one potential outcome; the other is a \textbf{counterfactual}. This is the
fundamental problem of causal inference.

For a binary outcome, the population causal risk difference is
\[
  \Delta_{\mathrm{causal}}
  :=\Pbb(\po{1}=1)-\Pbb(\po{0}=1)
  =\E[\po{1}]-\E[\po{0}].
\]
The difference visible in observational data is instead
\[
  \Delta_{\mathrm{obs}}
  :=\Pbb(\rv{y}=1\mid\rv{t}=1)
    -\Pbb(\rv{y}=1\mid\rv{t}=0).
\]
By consistency, the first term describes $\po{1}$ among treated units and the
second describes $\po{0}$ among control units. These are generally different
subpopulations, so association need not equal causation.

If treatment is independent of both potential outcomes,
\[
  \rv{t}\indep\{\po{0},\po{1}\},
\]
then $\Delta_{\mathrm{obs}}=\Delta_{\mathrm{causal}}$ at the population level.
Random assignment creates this independence by design. In a finite randomized
sample, the two observed rates are still estimates and need not equal the
population probabilities exactly.

\section{Confounding and adjustment}

A pre-treatment variable $\rv{a}$ is a \textbf{confounder} when it helps create
dependence between treatment and the potential outcomes. For example, age may
affect which patients take a drug and may also predict recovery under either
treatment. Comparing treated and control patients without accounting for age
then mixes together treatment differences and age-composition differences.

Suppose the following assumptions are plausible:
\begin{enumerate}[label=\arabic*.,leftmargin=2em,itemsep=0.2em]
  \item \textbf{Consistency:} the observed outcome equals the potential outcome
        under the treatment received.
  \item \textbf{Conditional independence:}
        $\{\po{0},\po{1}\}\indep\rv{t}\mid\rv{a}$.
  \item \textbf{Positivity:} within every relevant value of $a$, both treatment
        groups can occur.
\end{enumerate}
Then, for $r\in\{0,1\}$,
\[
  \boxed{
  \Pbb(\po{r}=1)
  =\sum_a \Pbb(\rv{a}=a)
     \Pbb(\rv{y}=1\mid\rv{a}=a,\rv{t}=r).}
\]
The same population weights $\Pbb(\rv{a}=a)$ are used for both treatments.
By contrast, the unadjusted observed rate is
\[
  \Pbb(\rv{y}=1\mid\rv{t}=r)
  =\sum_a \Pbb(\rv{a}=a\mid\rv{t}=r)
     \Pbb(\rv{y}=1\mid\rv{a}=a,\rv{t}=r),
\]
which may use very different weights in the two treatment groups. This is the
law of total probability, and the changing weights are the mechanism behind
Simpson's paradox.

\section{Simpson's paradox: Curry versus Lee}

During the 2014--2015 NBA season, Courtney Lee had the higher overall
three-point percentage, but Stephen Curry had the higher percentage within
both distance categories:
\begin{center}
\begin{tabular}{@{}lcc@{}}
\toprule
Shot type & Stephen Curry & Courtney Lee \\
\midrule
Short threes ($\leq 24$ feet) & $45/90=50.0\%$ & $56/116=48.3\%$ \\
Long threes ($>24$ feet) & $145/366=39.6\%$ & $19/55=34.5\%$ \\
\midrule
All threes & $190/456=41.7\%$ & $75/171=43.9\%$ \\
\bottomrule
\end{tabular}
\end{center}
This reversal is \textbf{Simpson's paradox}. Shot distance, denoted $\rv{d}$,
explains it. Lee took $116/171=67.8\%$ of his attempts from the easier short
range, while Curry took only $90/456=19.7\%$ from that range. Each overall
percentage is a weighted average with a different set of weights. For example,
\[
  \Pbb(\rv{y}=1\mid\rv{t}=\text{Lee})
  =0.678(0.483)+0.322(0.345)=0.439.
\]

For a common comparison, weight both players by the shot mix in the combined
data: $206/627=0.329$ short and $421/627=0.671$ long. The standardized rates are
\begin{align*}
  \text{Curry: }&0.329(0.500)+0.671(0.396)=0.430,\\
  \text{Lee: }&0.329(0.483)+0.671(0.345)=0.390.
\end{align*}
Standardization therefore favors Curry. It is a useful comparison of shooting
under a common distance distribution. Interpreting it as a causal comparison
requires the strong assumption that, conditional on distance, player identity
is independent of the relevant potential outcomes; other shot characteristics
may violate that assumption.

\keybox{\textbf{Read the comparison before interpreting it.}
Marginal rates answer ``what happened under each group's actual mix?''
Stratum-specific rates answer ``what happened within a fixed category?''
Standardized rates answer ``what would the rates be under a common mix?''}

\newpage
\section*{Practice problems}

\subsection*{Problem 1: What is observed? \hfill [7 minutes]}
Four patients have binary treatment and recovery outcomes:
\begin{center}
\begin{tabular}{@{}ccccc@{}}
\toprule
Patient & Treatment $t$ & Observed $y$ & $\mathrm{po}_0$ & $\mathrm{po}_1$ \\
\midrule
A & 0 & 0 & \phantom{0} & \phantom{0} \\
B & 0 & 1 & \phantom{0} & \phantom{0} \\
C & 1 & 1 & \phantom{0} & \phantom{0} \\
D & 1 & 0 & \phantom{0} & \phantom{0} \\
\bottomrule
\end{tabular}
\end{center}
\begin{enumerate}[label=(\alph*),itemsep=0.4em]
  \item Fill in every potential outcome determined by the observed data. Mark
        each unobserved potential outcome with ``?''.
  \item Compute the observed difference in recovery rates, treated minus control.
  \item Assuming only that potential outcomes are binary, find the smallest and
        largest possible finite-sample average causal effect
        $\frac14\sum_{i=1}^4(\mathrm{po}_{1,i}-\mathrm{po}_{0,i})$.
  \item If treatment was randomly assigned, does that reveal each missing
        counterfactual? What does randomization provide instead?
\end{enumerate}
\workingspace{5.0cm}

\begin{answer}
Consistency determines the entries under the treatment actually received:
\begin{center}
\begin{tabular}{@{}ccccc@{}}
\toprule
Patient & $t$ & $y$ & $\mathrm{po}_0$ & $\mathrm{po}_1$ \\
\midrule
A & 0 & 0 & 0 & ? \\
B & 0 & 1 & 1 & ? \\
C & 1 & 1 & ? & 1 \\
D & 1 & 0 & ? & 0 \\
\bottomrule
\end{tabular}
\end{center}
Both groups recover at rate $1/2$, so the observed difference is $0$.

The known $\mathrm{po}_1$ values sum to $1$, and the two missing values can add
anything from $0$ to $2$. Thus $\sum_i\mathrm{po}_{1,i}\in[1,3]$. The same
argument gives $\sum_i\mathrm{po}_{0,i}\in[1,3]$. The difference of the sums
can range from $-2$ to $2$, so the finite-sample average causal effect can range
from $-1/2$ to $1/2$. The observed difference $0$ does not determine it.

Randomization does not reveal any patient's missing counterfactual. It makes
treatment independent of potential outcomes, so group averages are unbiased
estimators of the corresponding population potential-outcome means (and become
accurate with sufficiently large samples).
\end{answer}

\newpage
\subsection*{Problem 2: Diagnose the Curry--Lee reversal \hfill [8 minutes]}
Use the basketball table in the review.
\begin{enumerate}[label=(\alph*),itemsep=0.45em]
  \item Which player has the higher marginal three-point percentage? Which
        player has the higher percentage within each distance category?
  \item Write Curry's overall percentage as a weighted average of his short- and
        long-shot percentages. Compare his weights with Lee's.
  \item In one or two sentences, explain why the direction reverses after the
        data are stratified by distance.
  \item Using a target shot mix of $40\%$ short and $60\%$ long, compute a
        standardized percentage for each player. Which is larger?
  \item Does part (d) by itself prove that Curry would outperform Lee on the
        same shots? State the additional assumption needed for that causal claim.
\end{enumerate}
\workingspace{5.5cm}

\begin{answer}
Lee has the higher marginal percentage: $43.9\%$ versus $41.7\%$. Curry has the
higher rate among both short attempts ($50.0\%$ versus $48.3\%$) and long
attempts ($39.6\%$ versus $34.5\%$).

Curry's overall rate is
\[
  \frac{90}{456}\frac{45}{90}
  +\frac{366}{456}\frac{145}{366}
  =0.197(0.500)+0.803(0.396)=0.417.
\]
Lee's weights are $116/171=0.678$ short and $55/171=0.322$ long. Lee's overall
rate receives much more weight from easier short attempts, which raises his
marginal percentage even though his rate is lower within each category.

With the common $40/60$ mix,
\begin{align*}
  \text{Curry: }&0.4(0.500)+0.6(0.396)=0.4376,\\
  \text{Lee: }&0.4(0.483)+0.6(0.345)=0.4002.
\end{align*}
Curry's standardized percentage is higher. A causal interpretation additionally
requires the relevant potential outcomes to be independent of player identity
conditional on distance. The table cannot establish this; other shot features
may still confound the comparison.
\end{answer}

\newpage
\subsection*{Problem 3: A drug and patient risk \hfill [8 minutes]}
An observational study records the following recoveries:
\begin{center}
\begin{tabular}{@{}lcc@{}}
\toprule
Patient group & Drug ($t=1$) & No drug ($t=0$) \\
\midrule
Lower risk & $45/50=90\%$ & $80/100=80\%$ \\
Higher risk & $30/100=30\%$ & $10/50=20\%$ \\
\midrule
All patients & $75/150=50\%$ & $90/150=60\%$ \\
\bottomrule
\end{tabular}
\end{center}
\begin{enumerate}[label=(\alph*),itemsep=0.45em]
  \item Compute the unadjusted observed risk difference. What conclusion would
        it suggest?
  \item Compute the risk difference within each risk group. Describe the reversal.
  \item In the full sample, half the patients are lower risk and half are higher
        risk. Standardize both treatment groups to this common $50/50$ mix.
        Compute the adjusted risk difference.
  \item Explain how risk group creates confounding here. State the assumptions
        needed to interpret the adjusted difference causally.
\end{enumerate}
\workingspace{5.5cm}

\begin{answer}
The unadjusted difference is
\[
  0.50-0.60=-0.10,
\]
which suggests that the drug lowers recovery by 10 percentage points. Within
the lower-risk group the difference is $0.90-0.80=0.10$, and within the
higher-risk group it is $0.30-0.20=0.10$. Thus the drug has the higher observed
recovery rate in each stratum even though it has the lower marginal rate.

Using the same $50/50$ weights for both treatment levels gives
\begin{align*}
  \Pbb(\po{1}=1)&\approx0.5(0.90)+0.5(0.30)=0.60,\\
  \Pbb(\po{0}=1)&\approx0.5(0.80)+0.5(0.20)=0.50.
\end{align*}
The adjusted risk difference is therefore $0.10$.

Only one third of the treated group is lower risk ($50/150$), compared with two
thirds of the untreated group ($100/150$). Risk group predicts recovery and is
distributed differently across treatment groups, so the marginal comparison is
confounded. The causal interpretation requires consistency, positivity, and
conditional independence of treatment and both potential outcomes given risk
group. It also assumes that the risk categories are measured well enough and
that no important confounder remains unmeasured.
\end{answer}

\ifsolutions\else\newpage\fi
\subsection*{Problem 4 (backup): True or false?}
\noindent\textit{Optional if time allows. Suggested working time: 5--7 minutes.}

\begin{enumerate}[label=(\alph*),itemsep=0.55em]
  \item Random assignment reveals both potential outcomes for every unit.
  \item If treatment has a positive risk difference in every stratum, then any
        standardized risk difference using the same nonnegative stratum weights
        for both treatments is positive.
  \item To adjust for a confounder, use the treated group's stratum weights for
        the treated rate and the control group's weights for the control rate.
  \item Seeing a Simpson reversal proves that one of the two marginal groups has
        been measured incorrectly.
\end{enumerate}
\workingspace{4.5cm}

\begin{answer}
The answers are false, true, false, false.
\begin{enumerate}[label=(\alph*),itemsep=0.35em]
  \item Randomization balances potential outcomes across treatment groups in
        distribution; it does not reveal a unit's counterfactual.
  \item A common-weight standardized difference is a weighted average of the
        positive stratum-specific differences, so it is positive.
  \item Those are exactly the treatment-specific weights used by the unadjusted
        marginal rates. Adjustment uses one common target distribution.
  \item The reversal can arise from different stratum compositions, even when
        every count is correct.
\end{enumerate}
\end{answer}

\section*{Takeaways}
\begin{itemize}[leftmargin=1.6em,itemsep=0.25em]
  \item Potential outcomes define causal questions, but only one is observed per unit.
  \item Randomization breaks dependence between treatment and potential outcomes.
  \item Confounding makes observed treatment groups incomparable; conditional
        comparisons and common-weight standardization can repair this only under
        substantive assumptions.
  \item Simpson's paradox is a weighted-average reversal, not an arithmetic error.
\end{itemize}

\section*{Sources}
The review follows the course handouts \emph{Causal Inference} and
\emph{Modeling Multiple Discrete Variables}, together with
\href{https://github.com/cfgranda/ps4ds/blob/main/slides/multiple%20discrete%20variables/simpsons_paradox_handout.pdf}{Was Courtney Lee Better Than Stephen Curry? Simpson's Paradox}
from \emph{Probability and Statistics for Data Science}.

\end{document}
